%&bericht

%%%%%%%%%%%%%%%%%%%%%%%%%%%%%%%%%%%%%%%%%%%%%%%%%%%%%%%%%%%%%%%%%%%%%%%%%%%%%%%
%% Descr:       Vorlage für Berichte der DHBW-Karlsruhe
%% Author:      Prof. Dr. Jürgen Vollmer, juergen.vollmer@dhbw-karlsruhe.de
%% $Id: bericht.tex,v 1.25 2020/03/13 15:07:45 vollmer Exp $
%%  -*- coding: utf-8 -*-
%%%%%%%%%%%%%%%%%%%%%%%%%%%%%%%%%%%%%%%%%%%%%%%%%%%%%%%%%%%%%%%%%%%%%%%%%%%%%%%

\documentclass[
   ngerman          % neue deutsche Rechtschreibung
  ,a4paper 
  ,oneside         % Papiergrösse
% ,twoside          % Zweiseitiger Druck (rechts/links)
% ,10pt             % Schriftgrösse
% ,11pt
  ,12pt
  ,pdftex
  ,parskip=half
  ,headsepline
  ,abstraction
  ,footsepline  
%  ,disable         % Todo-Markierungen auschalten
]{report}
\usepackage[ngerman]{babel}
\usepackage[T1]{fontenc}
% Bitte die Codierung Ihrer Dateien auswählen:
% \usepackage[latin1]{inputenc}    % Für UNIX mit ISO-LATIN-codierten Dateien
% \usepackage[applemac]{inputenc}  % Für Apple Mac
% \usepackage[ansinew]{inputenc}   % Für Microsoft Windows
\usepackage[utf8]{inputenc}        % UTF-8 codierte Dateien
                                   % Dieses Dokument ist unter Unix erstellt, daher
                                   % wird diese Input-Codierung benutzt.

\usepackage{bericht}


%% ACHTUNG, wenn man eine eigene Formatdatei (bericht.fmt) benutzt, werden Änderungen an bericht.sty
%% erst wirksam, wenn die Format-Datei neu erzeugt wurde!!!
%% Genauer alle Änderungen, die textuell vor der nächsten Zeile ".... endofdump...." stehen
%% werden erst wirksam, wenn die Formatdatei neu erzeugt wurde
\csname endofdump\endcsname

%%%%%%%%%%%%%%%%%%%%%%%%%%%%%%%%%%%%%%%%%%%%%%%%%%%%%%%%%%%%%%%%%%%%%%%%%%%%%%%
%% Angaben zur Arbeit
%%%%%%%%%%%%%%%%%%%%%%%%%%%%%%%%%%%%%%%%%%%%%%%%%%%%%%%%%%%%%%%%%%%%%%%%%%%%%%%

\newcommand{\Autor}{Matthias Schuhmacher und Marlene Rieder}
\newcommand{\MatrikelNummer}{4128647 und 8261867}
\newcommand{\Kursbezeichnung}{tinf21b3 und tinf21b5}

%\newcommand{\FirmenName}{Siemens AG}
%\newcommand{\FirmenStadt}{Karlsruhe}
%\newcommand{\FirmenLogoDeckblatt}{{\includegraphics[width=5cm]{Siemens_Logo_transp}}}

% Falls es kein Firmenlogo gibt:
%  \newcommand{\FirmenLogoDeckblatt}{}

%\newcommand{\BetreuerFirma}{Boris Rojs}
\newcommand{\BetreuerDHBW}{Prof. Dr. Marcus Strand}

%%%%%%%%%%%%%%%%%%%%%%%%%%%%%%%%%%%%%%%%%%%%%%%%%%%%%%%%%%%%%%%%%%%%%%%%%%%%%%%%%%%%%

% Wird auf dem Deckblatt und in der Erklärung benutzt:
%\newcommand{\Was}{Projekt-/Studien-/Bachleorarbeit}
%\newcommand{\Was}{Projektarbeit}
\newcommand{\Was}{Studienarbeit}
%\newcommand{\Was}{Bachleorarbeit}

%%%%%%%%%%%%%%%%%%%%%%%%%%%%%%%%%%%%%%%%%%%%%%%%%%%%%%%%%%%%%%%%%%%%%%%%%%%%%%%%%%%%%

\newcommand{\Titel}{Pepper VR – Teleopperation eines humanuiden Roboter auf Basis der Analyse menschlicher Bewegung}
\newcommand{\AbgabeDatum}{20. Mai 2024}

\newcommand{\Dauer}{300 Stunden}

% \newcommand{\Abschluss}{Bachelor of Engineering}
\newcommand{\Abschluss}{Bachelor of Science}

%\newcommand{\Studiengang}{Informatik / Informationstechnik}
 \newcommand{\Studiengang}{Informatik / Angewandte Informatik}

\hypersetup{%%
  pdfauthor={\Autor},
  pdftitle={\Titel},
  pdfsubject={\Was}
}

%%%%%%%%%%%%%%%%%%%%%%%%%%%%%%%%%%%%%%%%%%%%%%%%%%%%%%%%%%%%%%%%%%%%%%%%%%%%%%%

% Wenn \includeonly{..} benutzt wird, werden nur diese Kaptitel ausgegeben.
%\includeonly{
 % abk
%  Einleitung
 %,Grundlagen
 %,Die Projekte
 %,Reflexion und Ausblick
%,changelog%
%}

%%%%%%%%%%%%%%%%%%%%%%%%%%%%%%%%%%%%%%%%%%%%%%%%%%%%%%%%%%%%%%%%%%%%%%%%%%%%%%%

% Benutzt man das "biblatex"-Paket, dann muß das hier stehen:
% siehe auch die mit BIBLATEX markierten Zeilen in bericht.sty
\bibliography{T1000.bib}
\def\SymbReg{\textsuperscript{\textregistered}}
\nocite{*}
\begin{document}

%%%%%%%%%%%%%%%%%%%%%%%%%%%%%%%%%%%%%%%%%%%%%%%%%%%%%%%%%%%%%%%%%%%%%%%%%%%%%%%

\begin{titlepage}
\begin{center}
\vspace*{-2cm}
\includegraphics[width=4cm]{dhbw-logo}\\[2cm]
{\Huge \Titel}\\[1cm]
{\Huge\scshape \Was}\\[1cm]
{\large für die Prüfung zum}\\[0.5cm]
{\Large \Abschluss}\\[0.5cm]
{\large des Studienganges \Studiengang}\\[0.5cm]
{\large an der}\\[0.5cm]
{\large Dualen Hochschule Baden-Württemberg Karlsruhe}\\[0.5cm]
{\large von}\\[0.5cm]
{\large\bfseries \Autor}\\[1cm]
{\large Abgabedatum \AbgabeDatum}
\vfill
\end{center}
\begin{tabular}{l@{\hspace{2cm}}l}
Bearbeitungszeitraum	         & \Dauer 			\\
Matrikelnummer	                 & \MatrikelNummer		\\
Kurs			         & \Kursbezeichnung		\\
%Ausbildungsfirma	         & \FirmenName			\\
%			         & \FirmenStadt			\\
%Betreuer der Ausbildungsfirma	 & \BetreuerFirma		\\
Gutachter der Studienakademie	 & \BetreuerDHBW		\\
\end{tabular}
\end{titlepage}

%%%%%%%%%%%%%%%%%%%%%%%%%%%%%%%%%%%%%%%%%%%%%%%%%%%%%%%%%%%%%%%%%%%%%%%%%%%%%%%

%%%%%%%%%%%%%%%%%%%%%%%%%%%%%%%%%%%%%%%%%%%%%%%%%%%%%%%%%%%%%%%%%%%%%%%%%%%%%%%

%% Descr:       Vorlage für Berichte der DHBW-Karlsruhe, Erklärung

%% Author:      Prof. Dr. Jürgen Vollmer, vollmer@dhbw-karlsruhe.de

%% $Id: erklaerung.tex,v 1.11 2020/03/13 14:24:42 vollmer Exp $

%% -*- coding: utf-8 -*-

%%%%%%%%%%%%%%%%%%%%%%%%%%%%%%%%%%%%%%%%%%%%%%%%%%%%%%%%%%%%%%%%%%%%%%%%%%%%%%%

 

% In Bachelorarbeiten muss eine schriftliche Erklärung abgegeben werden.

% Hierin bestätigen die Studierenden, dass die Bachelorarbeit, etc.

% selbständig verfasst und sämtliche Quellen und Hilfsmittel angegeben sind. Diese Erklärung

% bildet das zweite Blatt der Arbeit. Der Text dieser Erklärung muss auf einer separaten Seite

% wie unten angegeben lauten.

 

\newpage
\thispagestyle{empty}
\begin{framed}
\begin{center}
\Large\bfseries Erklärung
\end{center}
\medskip
\noindent
% siehe §5(3) der \enquote{Studien- und Prüfungsordnung DHBW Technik} vom 29.\,9.\,2017 und Anhang 1.1.13
Ich versichere hiermit, dass ich meine \Was mit dem Thema:
\enquote{\Titel}
selbstständig verfasst und keine anderen als die angegebenen Quellen und Hilfsmittel benutzt habe. Ich versichere zudem, dass die eingereichte elektronische Fassung mit der gedruckten Fassung übereinstimmt.
\vspace{3cm}
\\
\noindent
\underline{\hspace{4cm}}\hfill\underline{\hspace{6cm}}\\
Ort~~~~~Datum\hfill Unterschrift\hspace{4cm}
\end{framed} 

\vfill
\begin{framed}
\begin{center}
\Large\bfseries Sperrvermerk
\end{center}
\medskip
\noindent
Die vorliegende Projektarbeit T3\_2000 mit dem Titel
\enquote{\Titel}
enthält unternehmensinterne bzw. vertrauliche Informationen der Atruvia AG, ist deshalb mit einem Sperrvermerk versehen und wird ausschließlich zu Prüfungszwecken am Studiengang \Studiengang der Dualen Hochschule Baden-Württemberg Karlsruhe vorgelegt.
\newline
Sie ist ausschließlich zur Einsicht durch den zugeteilten Gutachter, die Leitung des Studiengangs und ggf. den Prüfungsausschuss des Studiengangs bestimmt. Es ist untersagt,
\newline
\indent - den Inhalt dieser Arbeit (einschließlich Daten, Abbildungen, Tabellen, Zeichnungen
\indent   usw.) als Ganzes oder auszugsweise weiterzugeben. Kopien oder Abschriften dieser
\indent   Arbeit (einschließlich Daten, Abbildungen, Tabellen, Zeichnungen usw.) als Ganzes
\indent   oder auszugsweise anzufertigen.
    \newline    
\indent - diese Arbeit zu veröffentlichen bzw. digital, elektronisch oder virtuell zur Verfügung zu
\indent   stellen.
    \newline
Jede anderweitige Einsichtnahme und Veröffentlichung - auch von Teilen der Arbeit - bedarf der vorherigen Zustimmung durch den Verfasser und die Atruvia AG.

\end{framed}

%%%%%%%%%%%%%%%%%%%%%%%%%%%%%%%%%%%%%%%%%%%%%%%%%%%%%%%%%%%%%%%%%%%%%%%%%%%%%%%
\endinput
%%%%%%%%%%%%%%%%%%%%%%%%%%%%%%%%%%%%%%%%%%%%%%%%%%%%%%%%%%%%%%%%%%%%%%%%%%%%%%%


%%%%%%%%%%%%%%%%%%%%%%%%%%%%%%%%%%%%%%%%%%%%%%%%%%%%%%%%%%%%%%%%%%%%%%%%%%%%%%%

\begin{abstract}
Das Ziel der vorliegenden Projektarbeit ist es, die agile Projektmanagement Methode Scrum auf Projekte der betrieblichen Ausbildung, im Speziellen auf das Umfeld der betrieblichen Ausbildung der Siemens AG, anzupassen. Dazu wurden zwei Projekte, die mit unterschiedlichen Ansätzen in der Ausbildungsabteilung durchgeführt wurden, betrachtet, Problematiken identifiziert und Lösungsansätze entwickelt. Aus den Betrachtungen geht hervor, dass das Scrum Framework grundsätzlich die Anforderungen der betrieblichen Ausbildungsabteilung der Siemens AG erfüllt, allerdings einige Anpassungen notwendig sind, auf welche im Kapitel 4 genauer eingegangen wird. Diese Arbeit richtet sich sowohl an bereits erfahrene Nutzer des Scrum Frameworks, als auch an Menschen die dieses Framework genauer kennenlernen möchten. Besonders widmet sich diese Arbeit an die Ausbilder im betrieblichen Umfeld, die an der Einführung eines neuen Projektmanagementansatzes interessiert sind und von den Erkenntnissen vergangener Projekte profitieren möchten.
\newline
\newline

\noindent
The goal of this project work was to adapt the agile project management method Scrum to projects of the company training, in particular to the environment of the company training of Siemens AG. For this purpose, two projects that were carried out with different approaches in the training department were examined, problems identified and solutions developed. From the observations it emerges that the Scrum framework basically meets the requirements of the corporate training department of Siemens AG, but that some adaptations are necessary, which will be discussed in more detail in Chapter 4. This work addresses both already experienced users of the Scrum framework and people who would like to get to know this framework in more detail. In particular, this work is dedicated to trainers in the operational environment who are interested in introducing a new project management approach and who would like to benefit from the findings of past projects.


\end{abstract}

\newpage
\tableofcontents           % Inhaltsverzeichnis hier ausgeben
\listoffigures             % Liste der Abbildungen
%\listoftables              % Liste der Tabellen
%\lstlistoflistings         % Liste der Listings
%\listofequations           % Liste der Formeln

% Jetzt kommt der "eigentliche" Text
%%%%%%%%%%%%%%%%%%%%%%%%%%%%%%%%%%%%%%%%%%%%%%%%%%%%%%%%%%%%%%%%%%%%%%%%%%%%%%%
%% Descr:       Vorlage für Berichte der DHBW-Karlsruhe, Datei mit Abkürzungen
%% Author:      Prof. Dr. Jürgen Vollmer, vollmer@dhbw-karlsruhe.de
%% $Id: abk.tex,v 1.4 2017/10/06 14:02:03 vollmer Exp $
%% -*- coding: utf-8 -*-
%%%%%%%%%%%%%%%%%%%%%%%%%%%%%%%%%%%%%%%%%%%%%%%%%%%%%%%%%%%%%%%%%%%%%%%%%%%%%%%

\chapter*{Abkürzungsverzeichnis}                   % chapter*{..} -->   keine Nummer, kein "Kapitel"
						         % Nicht ins Inhaltsverzeichnis
% \addcontentsline{toc}{chapter}{Akürzungsverzeichnis}   % Damit das doch ins Inhaltsverzeichnis kommt

% Hier werden die Abkürzungen definiert
\begin{acronym}[DHBW]
  % \acro{Name}{Darstellung der Abkürzung}{Langform der Abkürzung}
  \acro{Abk}[Abk.]{Abkürzung}
  \acro{DHBW} {Duale Hochschule Baden-Württemberg}
  \acro{SAS}[SAS]{Société par actions simplifiée}
  \acro{B2B}[B2b]{Business-to-Business}
  \acro{TH}[TH]{Technische Hochschule}
  \acro{API}[API]{Application Programming Interface}
  \acro{SDK}[SDK]{Software Development Kit}
  \acro{URL}[URL]{Uniform Resource Locator}
  \acro{IP}[IP]{Internet Protocol}
  \acro{RPC}[RPC]{Remote Procedure Call}
  \acro{LPC}[LPC]{Local Procedure Call}
  
  % Folgendes benutzen, wenn der Plural einer Abk. benöigt wird
  % \newacroplural{Name}{Darstellung der Abkürzung}{Langform der Abkürzung}
  \newacroplural{Abk}[Abk-en]{Abkürzungen}
  % Wenn neicht benutzt, erscheint diese Abk. nicht in der Liste
  \acro{NUA}{Not Used Acronym}
\end{acronym}
              % Abkürzungsverzeichnis
\chapter{Einleitung}
Im Umfeld der \ac{DHBW} Karlsruhe werden immer wieder Projekte im Bereich der Robotik an Studierende im Rahmen verschiedener Arbeiten vergeben. Ebenso dieses Projekt, welches im Rahmen der Studienarbeit an die beiden Studierenden Matthias Schuhmacher und Marlene Rieder übergeben wurde. 
\\

\noindent
Im Umfeld der \ac{DHBW} Karlsruhe werden immer wieder Projekte im Bereich der Robotik an Studierende im Rahmen verschiedener Arbeiten vergeben. Ebenso dieses Projekt, welches im Rahmen der Studienarbeit an die beiden Studierenden Matthias Schuhmacher und Marlene Rieder übergeben wurde.
\\

\noindent 
In unserer immer digitaler werdenden Welt spiele sowohl Robotik als auch \ac{VR} eine immer größer werdende Rolle, diese zwei Welten sollen durch das, im folgenden beschriebene Projekt verbunden werden.
\\

\noindent
Ziel dieses Projektes ist es den humanoiden Roboter Pepper mit Hilfe der \ac{VR} Brille Meta Quest 3 zu steuern. Dies soll ermöglicht werden, indem das Kamerabild des Roboters an die Brille übertragen wird und sich die Bewegungen des Roboters durch die Controller der Brille steuern lassen. Der Nutzer soll also sehen, was sich vor dem Roboter befindet und ihn dann steuern können.

\section{Motivation}
Die Teleoperation humanoider Roboter bietet eine Vielzahl von Vorteilen. Durch die Fernsteuerung können Roboter in gefährlichen oder unzugänglichen Umgebungen eingesetzt werden, was die Sicherheit der menschlichen Bediener erhöht. Außerdem ermöglicht die Verwendung von VR-Technologie eine intuitivere und natürlichere Steuerung der Roboter, da die Bediener das Gefühl haben, direkt vor Ort zu sein und in Echtzeit mit der Umgebung zu interagieren.

\section{Ziele und Fragestellungen}
Das Hauptziel dieser Arbeit ist es, eine funktionsfähige Verbindung zwischen der MetaQuest 3 VR-Brille und dem humanoiden Roboter Pepper herzustellen. Dabei sollen folgende Fragestellungen untersucht werden:
\begin{itemize}
    \item Wie kann die Datenübertragung zwischen der VR-Brille und dem Roboter optimiert werden, um eine latenzfreie Steuerung zu ermöglichen?
    \item Welche Herausforderungen ergeben sich bei der Integration von Unity und ROS zur Steuerung des Roboters?
    \item Welche technischen und ergonomischen Anforderungen müssen bei der Entwicklung der VR-Steuerung berücksichtigt werden?
\end{itemize}

\section{Methodik}
Zur Beantwortung dieser Fragestellungen wird ein methodischer Ansatz verfolgt, der sowohl theoretische als auch praktische Aspekte umfasst. Zunächst wird eine umfassende Literaturrecherche durchgeführt, um den aktuellen Stand der Technik und die verfügbaren Technologien zu analysieren. Anschließend wird ein Prototyp entwickelt, der die Verbindung zwischen der VR-Brille und dem Roboter herstellt. Dieser Prototyp wird in mehreren Iterationen getestet und optimiert, um die bestmögliche Leistung zu erzielen.

\section{Struktur der Arbeit}
Diese Arbeit gliedert sich wie folgt:
\begin{itemize}
    \item Kapitel 2 beschreibt die theoretischen Grundlagen und den aktuellen Stand der Technik.
    \item Kapitel 3 erläutert die technischen Spezifikationen der verwendeten Hardware und Software.
    \item Kapitel 4 stellt den Entwicklungsprozess und die Implementierung des Prototyps dar.
    \item Kapitel 5 präsentiert die Testergebnisse und Diskussion.
    \item Kapitel 6 gibt eine Zusammenfassung und Ausblick auf zukünftige Arbeiten.
\end{itemize}


\chapter{Grundlagen}\label{chap:Grundlagen}

\section{Pepper}\label{sec:Pepper}
Pepper ist ein humanoider Roboter, der entwickelt wurde, um die Gefühle und Gesten von Menschen zu analysieren und basierend auf diesen, darauf zu reagieren. Das Projekt entstand durch eine Zusammenarbeit des französischen Unternehmens Aldebaran Robotics \ac{SAS} und des japanischen Telekommunikations- und Medienkonzerns SoftBank Mobile Corp. Ziel diese Projektes war es, einen humanuiden ``Roboter-Gefährten'' oder einen ``persönlichen Roboter-Freund'' zu schaffen, der zunächst im Gewerbesektor in Verkaufsräumen, an Empfangstischen oder in Bildungs- und Gesundheitseinrichtungen eingesetzt werden sollte. Die Produktion wurde jedoch aufgrund geringer Nachfrage bis auf Weiteres pausiert.\\
Das Konzept von Pepper distanziert sich von herkömmlichen Industrierobotern und reinen Spielzeugrobotern, indem er als informativer und kommunikativer Begleiter konzipiert wurde. Sein Aussehen, das im etwa an die Größe eines Kindes angelehnt ist, sowie ein freundliches Gesicht und eine kindliche Stimme sind im ästhetischen Konzept von ``kawaii'' (japanisch für ``niedlich'' oder auch ``liebenswert'') gehalten.\\
Pepper wurde im Rahmen einer Präsentation am 5. Juni 2014 als der ``erste persönliche Roboter der Welt mit Emotionen'' vorgestellt. Die Vermarktung begann damit, dass SoftBank Pepper-Geräte in ihren Verkaufsräumen einsetzte, um Kunden zu unterhalten und zu informieren. Die Roboter sollte dabei den Umgang mit Kunden erlernen, um zukünftige Anwendungsmöglichkeiten zu erforschen. Verkauft wurde offiziell ab dem 3. Juli 2015 zu einem Preis von 198.000 Yen pro Einheit, zuzüglich monatlicher Gebühren für Zusatzleistungen. Im Laufe der Zeit wurde Pepper auch für den Einsatz in weiteren Unternehmen und Einrichtungen verfügbar gemacht.\\
Pepper wird mit einer Grundausstattung an Anwendungen geliefert, jedoch sind für spezifische Anwendungen, individuell entwickelte Softwarelösungen erforderlich wie auch zum Beispiel in diesem. SoftBank ermöglichte unabhängigen Entwicklern durch die Veröffentlichung der Schnittstellen den Zugang zu einem Interface für Applikationsprogramme, um zusätzliche Anwendungen für Pepper zu erstellen. Das NAOqi-Framework welches für diesen Nutzen bereitgestellt wurde, beinhaltet eine \ac{API}, eine \ac{SDK} und weitere Tools, welche in den Sprachen Python und C++ uneingeschränkten Zugriff auf die Komponenten, Sensoren und Aktoren des Roboters bieten, dazu später Ausführlicheres in \autoref{sec:NAOqi}. Mit Hilfe diese Interfaces haben verschiedene Unternehmen integrierte Lösungen entwickelt, die Pepper beispielsweise bei der Kundenberatung unterstützen können.\\
Das Design von Pepper ist dem Menschen ähnlich und umfasst einen Kopf mit integrierten Mikrofonen und Kameras sowie einen Torso mit weiteren Sensoren für Stabilität und Sicherheit. Der Roboter verfügt über verschiedene Mechaniken, die es ihm ermöglichen, sich flüssig zu bewegen und mit Personen zu interagieren. Durch die Verwendung von Kameras und bereitgestellter Software ist Pepper in der Lage, Emotionen bei seinen Gesprächspartnern zu erkennen und darauf zu reagieren, obwohl er selbst keine Mimik besitzt. Sicherheitsvorkehrungen wie Abstandssensoren und Stabilisatoren gewährleisten einen sicheren Einsatz von Pepper in verschiedenen Umgebungen. Diese können jedoch bedingt durch den Entwickler deaktiviert werden, um den Roboter in komplexeren oder laborähnlichen Umgebungen zu betreiben.\\

\section{Publish and Subscribe}\label{sec:PublishAndSubscribe}
Bevor wir uns mit den Technologien und Herangehensweisen des Pepper Roboters beziehungsweise dessen Betriebssystem beschäftigen, müssen wir uns mit dem Publish-Subscribe-Modell auseinandersetzen.\\
Das Publish-Subscribe-Modell ist ein Paradigma für einen effektiven Nachrichtenaustausch in verteilten Systemen. Erstmals publiziert in einem Paper der \ac{ACM} von 1987\cite{wiki_publish_subscribe_pattern}, ermöglicht es die flexible Kommunikation zwischen verschiedenen Komponenten, indem es einen Mechanismus bereitstellt, über den Nachrichten von einem Sender, dem Publisher, an einen oder mehrere Empfänger, den Subscribern, verteilt werden können. Ein zentrales Element dieses Modells ist dabei Nachrichtenbroker oftmals auch nur als Borker referenziert, der als Vermittler zwischen Publishern und Subscribenden fungiert.\\
Der Broker empfängt Nachrichten vom Publisher und organisiert sie in verschiedene Kategorien, welche als Topics bezeichnet werden. Diese Topics dienen als thematische Selektionsmerkmale, die den Inhalt der Nachrichten beschreiben. Sie können in den verschiedensten Datentypen erfolgen, einfache Ganzzahlen, Texte bis hin zu Bildern oder weitaus komplexeren Datenstrukturen. Durch diese Kategorisierung in den Topics, wird eine gezielte Auswahl und Weiterleitung der Nachrichten ermöglicht.\\
Subscriber können bestimmte Topics abonnieren, die ihren Gewünschten entsprechen. Dadurch erhalten sie nur die Nachrichten, die zu den von ihnen gewählten Themen gehören. Dieser Prozess der Nachrichtenauswahl und -verarbeitung wird als Filterung bezeichnet und kann auf zwei Arten erfolgen: themenbasiert und inhaltsbasiert.\\
Im themenbasierten Ansatz erhalten Abonnenten alle Nachrichten zu den Topics, welche sie abonniert haben. Der Publisher definiert die verfügbaren Topics, aus welchen die Abonnenten selbst wählen können. Im inhaltsbasierten Ansatz dagegen, werden Nachrichten nur an die Abonnenten weitergeleitet, wenn sie den vom dessen festgelegten Kriterien entsprechen. Dabei ist der Abonnent für die Spezifikation dieser Kriterien verantwortlich\cite{itwissen_publish_subscribe_model}.\\
Das Publish-Subscribe-Modell bietetdurch seine Herangehensweise eine hohe Flexibilität und Skalierbarkeit, und wird daher in verschiedenen Anwendungsbereichen eingesetzt. Es ermöglicht eine Entkopplung von Nachrichtenerzeugung und -verarbeitung, was gerade in der Entwicklung verteilter Systeme Prozesse erleichtert. Die vermutlich bekannteste Implementierung des Modell ist das MQTT-Protokoll, welches konzipiert wurde um Telemetriedeaten zwischen Sensoren und Servern zu übertragen speziell in unzuverlässigen Umgebungen\cite{elektronik_kompendium_publish_subscribe}.\\
Weiter findet das Publish-Subscribe-Modell Anwendung in einer Vielzahl von weiteren Systemen und Technologien. Beispiele dafür sind der \ac{JMS}, der \ac{DDS} oder wie im Fall diese Projektes die \ac{MOM} zu denen auch \ac{ROS} gehört.\\

\section{NAOqi}\label{sec:NAOqi}
Im folgenden Abschnitt wird das NAOqi-Framework, welches auf dem Pepper Roboter läuft, genauer erläutert.\\

\subsection{Was ist NAOqi?}\label{subsec:WasIstNAOqi}
NAOqi ist die Bezeichnung für die Hauptsoftware, die auf dem Pepper Roboter ausgeführt wird und ihn intern steuert. Diese kann mit persönlichen Modulen weiterentwickelt und angepasst werden. Dazu können mit Hilfe der Aldebaran \ac{SDK} auch NQOqi-\ac{SDK} genannt, eigene Module oder Bibliotheken entwickelt werden. Diese NAOqi-\ac{SDK} ist die Basis des NAOqi Frameworks und ist in C++ implementiert\cite{nao_dev_install_guide}.\\
Das NAOqi Framework ist das Programmiergerüst, welches zur Programmierung von NAO und Pepper Robotern verwendet wird. Es implementiert alle allgemeinen Anforderungen der Robotik, einschließlich: Parallelität, Ressourcen-Management, Synchronisation und Ereignisse. Dieses Framework ermöglicht eine homogene Kommunikation zwischen verschiedenen Modulen wie etwa die Bewegung, Audio oder Video sowie eine homogene Programmierung und einen homogenen Informationsaustausch. Das Framework ist:\\
\begin{itemize}
    \item plattformübergreifend, wie bereits erwähnt, basiert die NAOqi-\ac{SDK} auf C++ und bietet damit die Möglichkeit auf Windows, Linux oder sogar auch Mac zu entwickeln.
    \item sprachübergreifend, mit einer identischen API für C++ und Python. Weitere Details dazu sind in \autoref{subsubsec:Sprachuebergreifend} aufgeführt.
    \item fähig auf Introspektion, was bedeutet, dass das Framework weiß, welche Funktionen in den verschiedenen Modulen verfügbar sind und wo. Für Details diesbezüglich siehe \autoref{subsubsec:Introspektion}.
\end{itemize}

\subsubsection{Sprachübergreifend}\label{subsubsec:Sprachuebergreifend}
Software kann in C++ und Python entwickelt werden. Eine Übersicht über die Sprachen selbst in den Abschnitten \autoref{subsec:Cpp} und \autoref{subsec:Python}. In allen Fällen sind die Programmiermethoden genau die gleichen, alle vorhandenen \acp{API} können unabhängig von den unterstützten Sprachen aufgerufen werden:\\
\begin{itemize}
    \item Wird ein neues C++-Modul erstellt, können die C++-\ac{API}-Funktionen von überall aus aufgerufen werden,
    \item Sind sie richtig definiert, können auch die \ac{API}-Funktionen eines Python-Moduls von überall aus aufgerufen werden.
\end{itemize}
Normalerweise werden die Verhaltensweisen in Python und Ihre Dienste in C++ entwickelt.\\
\normfig{NAOqi-Framework-Sprachübergreifend}{0.5}{-}

\subsubsection{Introspektion}\label{subsubsec:Introspektion}
Die Introspektion ist die Grundlage der Roboter-\ac{API}, der Fähigkeiten, der Überwachung und der Maßnahmen bei überwachten Funktionen. Der Roboter selbst kennt alle verfügbaren \ac{API}-Funktionen. Wird eine Bibliothek entladen, werden die entsprechenden \ac{API}-Funktionen automatisch ebenfalls entfernt. Eine in einem Modul definierte Funktion kann der \ac{API} mit einem \texttt{BIND\_METHOD} hinzugefügt werden.\\
Wird eine Funktion gebunden, werder automatisch folgende Funktionen ausgeführt:\\
\begin{itemize}
    \item Funktionsaufruf in C++ und Python, wie in \autoref{subsubsec:Sprachuebergreifend} beschrieben
    \item Erkennen der Funktion, wenn sie gerade ausgeführt wird
    \item Funktion lokal oder aus der Ferne, z.B. von einem Computer oder einem anderen Roboter, ausführen weiter im Detail beschrieben in \autoref{subsubsec:VerteilterBaumUndKommunikation}
    \item Generierung und Aufruf von \texttt{wait}, \texttt{stop}, \texttt{isRunning} in Funktionen
\end{itemize}
Die \ac{API} wird im Webbrowser angezeigt wenn auf das Gerät per \ac{URL} oder \ac{IP}-Addresse auf dem Port 9559 zugegriffen wird. In dieser Übersicht, zeigt der Roboter seine Modulliste, Methodenliste, Methodenparameter, Beschreibungen und Beispiele an. Der Browser zeigt auch parallele Methoden an, die überwacht, zum Warten veranlasst und gestoppt werden können.\\
Die Introspektion und derer Implementation im NAOqi-Framework, ist also ein nützliches Werkzeug, welches es ermöglicht, die Roboter-\ac{API} zu verstehen und zu verwenden aber auch zu überwachenund zu steuern.\\
\normfig{NAOqi-Framework-Introspektion}{0.5}{-}

\subsubsection{Verteilter Baum und Kommunikation}\label{subsubsec:VerteilterBaumUndKommunikation}
Eine Echtzeitanwendung kann aus einer einzelnen ausführbaren Datei oder einem Baum von mehreren Systemen wie etwa Robotern, Prozessen oder Modulen bestehen. Unabhängig davon sind die Aufrufmethoden immer dieselben. Eine ausführbare Datei kann durch eine Verbindung mit einem anderen Roboter mit \ac{IP}-Adresse und Port verbunden werden, sodass alle \ac{API}-Methoden von anderen ausführbaren Dateien sind auf die gleiche Weise verfügbar sind, genau wie bei einer lokalen Methode. NAOqi trifft dabei selbst die Wahl zwischen schnellem Direktaufruf \ac{LPC} und Fernaufruf \ac{RPC}.\\
\normfig{NAOqi-Framework-Kommunikation}{0.5}{-}

\subsection{NAOqi Vorgehensweise}\label{subsec:NAOqiVorgehensweise}
Die NAOqi Software, welche auf dem Roboter läuft, ist ein Broker. Wenn dieser startet, lädt er eine Voreinstellungsdatei in den Speicher, in der festgelegt ist, welche Bibliotheken in dieser Konfiguration geladen werden sollen. Jede Bibliothek enthält ein oder mehrere Module, die den Broker benutzen, um ihre Methoden bereitzustellen.\\
\normfig{NAOqi-Broker-Bibliotheken-Module}{0.5}{-}
Der Broker selbst bietet Nachschlagdienste an, so dass jedes Modul im Baum oder im Netzwerk jede Methode finden kann, die an dem Broker bekannt gegeben wurde.\\
Das Laden von Modulen bildet dann einen Baum von Methoden, die mit Modulen verbunden sind, und von Modulen, welche mit dem Broker verbunden sind.\\
\normfig{NAOqi-Broker-Modul-Methoden}{0.5}{-}

\subsubsection{NAOqi Proxy}\label{subsubsec:NAOqiProxy}
Ein weiterer wichtiger Bestandteil des Broker ist der Proxy welcher grundlegend die Aufgabe eines Modules repräsentiert. Dies kann lokal oder entfernt enstehen der einzige Unterschied dabei ist, dass bei einer entfernten Refernz die \ac{IP} des entfernten Broker mit angegeben werden muss, gleich bleibt aber im internen, dass der Proxy die Methoden des Moduls an den Broker weiterleitet.\\

\subsection{NAOqi Module}\label{subsubsec:NAOqiModule}

\section{Robot Operating System}\label{sec:ROS}

\section{Programming Languages}

\subsection{C++}\label{subsec:Cpp}

\subsection{Python}\label{subsec:Python}

\section{Virtual Reality}

\section{Entwicklung für Virtual Reality}
\section{TCP}
\include{DieProjekte}
\include{ReflexionUndAusblick}



% Ab hier beginnt der Anhang
\appendix
\addcontentsline{toc}{chapter}{Literaturverzeichnis}

% Haben Sie das "biblatex"-Paket nicht installiert, benutzen Sie folgendes:
% Ohne das "biblatex"-Paket (s. bericht.sty) produziert folgendes
% "deutsche" Zitate in Literaturverzeichnissen gemaß der Norm DIN 1505,
% Teil 2 vom Jan. 1984.
% Die Zitatmarken werden alphabetisch nach Verfassern
% sortiert und sind durch abgekürzte Verfasserbuchstaben plus
% Erscheinungsjahr in eckigen Klammern gekennzeichnet.

% \bibliographystyle{alphadin}
% \bibliography{bericht}

%%%%%%%%%%%%%%%%%%%%%%%%%%%%%%%%%%%%%%%5
% BIBLATEX
% Benutzt man das "biblatex"-Paket, muß man folgendes schreiben:
\def\refname{Literaturverzeichnis}
\printbibliography
%%%%%%%%%%%%%%%%%%%%%%%%%%%%%%%%%%%%%%%5


\end{document}
